\section{Conclusion}
\label{sec:conclusion}

We have presented a controlled study of prompt saturation, the point at which
response quality stops improving as prompt length increases.
Using a seven-level additive prompt methodology across six tasks and seven LLMs,
we find a consistent \textbf{structured--open dichotomy}: classification
saturates at 42--64 tokens, product extraction at 92--536 tokens, while QA and
math reasoning show no statistically significant saturation at all.

These findings support a nuanced view of prompt engineering: more tokens help
\emph{up to a task-specific threshold} for tasks requiring schema compliance,
and \emph{not at all} for tasks bottlenecked by model knowledge or reasoning.
A larger-scale replication on 200 examples from SST-2 and SQuAD confirms that
these saturation patterns persist when evaluated on substantially larger
datasets.
The practical implication is direct: prompt budgets should be set by task
type, not by general verbosity heuristics.

\paragraph{Future directions.}
Extending the study to additional structured tasks (code generation with type
annotations, data-to-text generation) would test whether the dichotomy
generalises.
Mechanistic interpretability analyses of how models process elaboration levels
could directly test the schema-compliance hypothesis.
Finally, studying saturation in multi-turn settings, where context accumulates
across turns, would address the increasingly common deployment pattern of
long-horizon LLM agents.
